\section{Vorläufige Gliederung}
\label{sec:vorlaufige-gliederung}


Vorläufige Gliederung

    1 Einleitung
    2 Kommunikation auf Twitter
        2.1 Funktionen von Twitter
        2.2 Hauptmerkmale von Twitter
        2.3 Politische Kommunikation auf Twitter
    3 Opinion-Leader und Two-Step-Flow
        3.1 Theoretischer Rahmen
        3.2 Theoretische Kritik am Two-Step-Flow
        3.3 Soziale Netzwerke als Schauplatz für Opinion-Leadership
    4 Methodik
    5 Donald Trump auf Twitter
        5.1 Trumps Nutzung von Twitter
        5.2 Case-Study der Kaufhauskette Nordstrom
    6 Diskussion und Fazit

Literaturverzeichnis

    Bennett, W. L., \& Manheim, J. B. (2006). The One-Step Flow of Communication. The ANNALS of the American Academy of Political and Social Science, 608(1), pp. 213–232.
    Crockett, Z. (2016). What I learned analysing 7 months of Donald Trump’s tweets. Vox. Retrieved from https://www.vox.com/2016/5/16/11603854/donald-trump-twitter
    Katz, E. (1957). The two-step flow of communication: An up-to-date report on an hypothesis. Public Opinion Quarterly, 21(1), 61–78.
    Katz, E., \& Lazarsfeld, P. (1955). Personal influence: The part played by people in the flow of mass communications. New York: Free Press.
    Lazarsfeld, L., Berelson, B., \& Gaudet, H. (1948). The people’s choice: How the voter makes up his mind in a presidential campaign. New York: Columbia University Press.
    Ott, B. (2017). The age of Twitter: Donald J. Trump and the politics of debasement. Critical Studies in Media Communication, 34(1), 59–68.
    Thorson, K. \& Wells, C. (2015). Curated flows: A framework for mapping media exposure in the digital age. Communication Theory, 26(3), 309–328